\documentclass[18pt]{article}
\usepackage{amsmath,amsfonts,amssymb}
\usepackage{natbib}
\usepackage{booktabs}
\usepackage[utf8]{inputenc}
\usepackage[T1]{fontenc}
\usepackage{amssymb}
\usepackage[version=4]{mhchem}
\usepackage{stmaryrd}
\usepackage{bbold}
\usepackage{nopageno}
\usepackage{hyperref}
\usepackage{bookmark}
\author{Adam Handwerger}
\pagestyle{plain}

\begin{document}
\begin{flushleft}
1 Orthogonal projection in a Hilbert space\\

\vspace{.4cm}
Let H be a Hilbert space and V a closed subspace of H. Note $\pi_V(f)$ is the orthogonal projection of f onto $V \in H$. Then \\
\vspace{.4cm}
1. $\pi_V(f) \in V$;\\
2. $\langle f-\pi_V(f),g\rangle=0 \text{ for all } g \in V$\\
\vspace{.4cm}
Show\\
\vspace{.4cm}
1. For any $f \in H, \parallel f \parallel^2=\parallel \pi_V(f) \parallel^2+\parallel f-\pi_V(f) \parallel^2$\\
\vspace{.4cm}
Since H is a Hilbert space, $f \in H$ can be decomposed as $f=f_V+f_{\perp}$.\\
Then\\
\vspace{.4cm}
$\parallel f \parallel^2=\langle f,f \rangle=\langle f_V+f_{\perp},f_V+f_{\perp} \rangle=$\\
$\langle f_V,f_V \rangle+2\langle f_V,f_{\perp} \rangle+\langle f_{\perp},f_{\perp} \rangle=$\\
$\parallel f_V \parallel^2+\parallel f_V-\pi_V(f) \parallel^2$\\
\vspace{.4cm}
This follows since $f_{\perp}$ is the orthogonal complement of $f_V$.\\
\vspace{.4cm}
2. For any $f,g \in H, \langle f,\pi_V(g)\rangle=\langle g,\pi_V(f)\rangle$\\
\vspace{.4cm}
$\langle f,\pi_V(g) \rangle=\langle f_V+f_{\perp},\pi_V(g) \rangle=\langle f_V,\pi_V(g) \rangle+\langle f_{\perp},\pi_V(g) \rangle=$\\
$\langle \pi_V(f),\pi_V(g) \rangle$ since $f_V =\pi_V(f)$\\
\vspace{.4cm}
In the same manner for g gives\\
\vspace{.4cm}
$\langle g,\pi_V(f) \rangle=\langle \pi_V(g),\pi_V(f) \rangle$\\
\vspace{.4cm}
Since the inner product is symmetric we can conclude that\\
\vspace{.4cm}
$\langle f,\pi_V(g) \rangle=\langle g,\pi_V(f) \rangle$\\
\vspace{.4cm}
3. Find an H and a subspace $V\subset H$ which is not closed in H and a $f\in H$ such that $\pi_V(f)$ does not exist.\\
\vspace{.4cm}
Let $H=L^2[-1,1]$ and V be the the subspace of all continuous functions f on $[-1,1]$\\
\vspace{.4cm}
Define a function $f^{\prime}$ as\\
$$
f^{\prime}(x)=
\begin{cases}
0, x \in [-1,-\frac{1}{2})\\
\frac{1}{2}, x \in [-\frac{1}{2},\frac{1}{2}]\\
0, x \in (\frac{1}{2},1]\\
\end{cases}
$$
\vspace{.4cm}
Note that $f^{\prime}$ is the limit of a sequence of continuous functions, but is not\\
continuous and there for is not contained in V. Therefore V is not closed.\\
Now consider the orthogonal complement of $f^{\prime}$, $f_{\perp} \in H$\\
\vspace{.4cm}
$$
f_{\perp}=
\begin{cases}
    \frac{1}{2}, x \in [-1,-\frac{1}{2})\\
    0, x \in [-\frac{1}{2},\frac{1}{2}]\\
    \frac{1}{2}, x \in (\frac{1}{2},1]\\
\end{cases}
$$

\vspace{.4cm}
$f_{\perp} \in H$, but $f^{\prime}$ is not in V. Therefore for $f=f^{\prime}+f_{\perp} \text{ , }\pi_V(f)$ does not\\
exist.\\
\vspace{.4cm}
2 A reproducing kernel is positive definite\\
\vspace{.4cm}
Let H be RKHS and $f:\mathcal{X} \mapsto \mathbb{R}$ by the kernel, K. This means\\
\vspace{.4cm}
1. for any $y \in \mathcal{X}, K(\cdot,y) \in H$\\
2. for any $f\in H$, and $y \in \mathcal{X}, \langle f,K(\cdot,y)\rangle=f(y)$\\
\vspace{.4cm}
1. Show that K is symmetric.\\
\vspace{.4cm}
Since $K \in H$ choose $f=K(\cdot,y^{\prime})$. Then by property 2.\\
$\langle K(\cdot,y^{\prime}),K(\cdot,y)\rangle=K(y^{\prime},y)$\\
\vspace{.4cm}
Likewise we also have\\
$\langle K(\cdot,y),K(\cdot,y^{\prime})\rangle=K(y,y^{\prime})$\\
\vspace{.4cm}
Since the inner product is symmetric $K(y^{\prime},y)=K(y,y^{\prime})$\\
\vspace{.4cm}
2. Show that K is positive definite.
\vspace{.4cm}
Let $x_i$ and $x_j \in \mathcal{X}$. Then for $\alpha_i \in \mathbb{R}$
$\sum\limits_i \sum\limits_j \alpha_i \alpha_j K(x_i,x_j)=\sum\limits_i \sum\limits_j \alpha_i \alpha_j \langle K(x_i,\cdot), K(x_j,\cdot) \rangle=$\\
$\langle \sum\limits_i \alpha_i K(x_i,\cdot),\sum\limits_j \alpha_j K(x_j,\cdot)\rangle=\parallel  \sum\limits_i \alpha_i K(x_i,\cdot) \parallel^2 \geq 0$\\
\vspace{.4cm}
3. Show for any $f\in H$ and $x \in \mathcal{X}, \lvert f(x) \rvert \leq \parallel f \parallel \sqrt{K(x,x)}$\\
\vspace{.4cm}
By the reproducing property of H, $f(x)=\langle f, K(\cdot,x) \rangle$. Then by the Cauchy-Schwarz inequality\\
\vspace{.4cm}
$\lvert f(x) \rvert^2=\lvert  \langle f, K(\cdot,x) \rangle \rvert^2 \leq \parallel f \parallel^2 \parallel K(\cdot,x) \parallel^2=$\\
$\parallel f \parallel^2 \langle K(\cdot,x),K(\cdot,x)\rangle=\parallel f \parallel^2 K(x,x) \implies$\\
$\lvert f(x) \rvert \leq \parallel f \parallel \sqrt{K(x,x)}$\\












\end{flushleft}
\end{document}